\vspace{-0.4cm}
\section{LLM Agents in Trust Games}
\vspace{-0.15cm}


\subsection{Trust Games}


Trust Games, referring to the Trust Game and its variations, have been widely used for examining human trust behavior in behavioral economics~\citep{berg1995trust,lentonIncentivisingTrust2011,Measuring_Trust,cesariniHeritabilityCooperativeBehavior2008}.
As shown in Figure~\ref{fig:framework}, the player who makes the first decision to send money is called the \textit{trustor}, while the other one who responds by returning money is called the \textit{trustee}. 
In this paper, we mainly focus on the following six types of Trust Games (the specific prompt for each game is articulated in the Appendix~\ref{Game Setting Prompt}):









\vspace{-3mm}
\paragraph{Game 1: Trust Game}
\label{game:Trustee trust rate Game}

 As shown in Figure~\ref{fig:framework}, in the Trust Game~\citep{coxHowIdentifyTrust2004b,berg1995trust}, the trustor initially receives $\$10$. The trustor selects $\$N$ and sends it to the trustee, exhibiting  \textit{trust behavior}. Then the trustee will receive $\$3N$, and have the option to return part of that $\$3N$ to the trustor, showing \textit{reciprocation behavior}. 

\vspace{-3mm}
\paragraph{Game 2: Dictator Game}
\label{game:Dictator Game and trust game}
In the Dictator Game~\citep{coxHowIdentifyTrust2004b}, the trustor also needs to send $\$N$  from the initial $\$10$ to the trustee and then the trustee will receive $\$3N$. Compared to the Trust Game, the only difference is that the trustee does not have the option to return money in the Dictator Game and the trustor is also aware that the trustee cannot reciprocate.


\vspace{-3mm}
\paragraph{Game 3: MAP Trust Game}
\label{game:MAP Problem}
In the MAP Trust Game (MAP represents Minimum Acceptable Probabilities)~\citep{BOHNET2004467}, a variant of the Trust Game,
the trustor needs to choose whether to trust the trustee. If the trustor chooses not to trust the trustee, each will receive $\$10$; If the trustor and the trustee both choose to trust, each will receive $\$15$; If the trustor chooses to trust, but the trustee does not, the trustor will receive $\$8$ and the trustee will receive $\$22$. There is probability $p$ that the trustee will choose to trust and $(1-p)$ probability that they will not choose to trust. MAP is defined as the minimum value of \( p \)  at which the trustor would choose to trust the trustee.




\vspace{-3mm}

\paragraph{Game 4: Risky Dictator Game} 
The Risky Dictator Game~\citep{BOHNET2004467} differs from the MAP Trust Game in only a single aspect. In the Risky Dictator Game,  the trustee is present but does not have the choice to trust or not and the money distribution relies on the pure probability $p$. Specifically, if the trustor chooses to trust, there is probability $p$ that both the trustor and the other player will receive $\$15$ and probability $(1-p)$ that the trustor will receive $\$8$ and the other player will receive $\$22$. If the trustor chooses not to trust the trustee, each player will receive $\$10$.







\vspace{-3mm}
\paragraph{Game 5: Lottery Game}
\label{game:Lottery Problem}
There are two typical Lottery Games~\citep{fetchenhauerBetrayalAversionPrincipled2012}. In the Lottery People Game, the trustor is informed that the trustee chooses to trust with probability \( p \). Then the trustor must choose between receiving fixed money or trusting the trustee, which is similar to the MAP Trust Game. In the Lottery Gamble Game, the trustor chooses between playing a gamble with a winning probability of \( p \) or receiving fixed money. \( p \) is set as \( 46\%\) following the human study.

\vspace{-3mm}
\paragraph{Game 6: Repeated Trust Game} We follow the setting of the Repeated Trust Game in~\citep{cochard2004trusting}, where the Trust Game is played for multiple rounds with the same players and each round begins anew with the trustor allocated the same initial money.



\vspace{-0.15cm}
\subsection{LLM Agent Setting}
\vspace{-0.15cm}
In our study, we set up our experiments using the CAMEL framework~\citep{li2023camel} with both closed-source and open-source LLMs including GPT-4, GPT-3.5-turbo-0613, GPT-3.5-turbo-16k-0613, text-davinci-003, GPT-3.5-turbo-instruct, Llama2-7b (or 13b, 70b) and Vicuna-v1.3-7b (or 13b, 33b)~\citep{instructgpt,openai2023gpt4,touvron2023Llama,vicuna2023}. We set the temperature as $1$ to increase the diversity of agents' decision-making and note that high temperatures are commonly adopted in related literature~\citep{aher2023using,lore2023strategic,guo2023gpt}. 

\mysection{Agent Persona} To better reflect the setting of real-world human studies~\citep{berg1995trust}, we design LLM agents with diverse personas in the prompt.  Specifically, we ask GPT-4 to generate 53 types of personas based on a given template. Each persona needs to have information including name, age, gender, address, job and background. Examples of the personas are shown in Appendix~\ref{Persona Prompt}.





\mysection{Belief-Desire-Intention (BDI)} 
The BDI framework is a well-established approach in agent-oriented programming~\citep{raoBDIAgentsTheory1995a} and was recently adopted to language models~\citep{andreas-2022-language-BDI}. We propose modeling LLM agents in Trust Games with the BDI framework to gain deeper insights into LLM agents' behaviors. Specifically, we let LLM agents directly output their Beliefs, Desires, and Intentions as the reasoning process for decision-making in Trust Games.































\begin{wrapfigure}{t}{0.57\textwidth}
\vspace{-4.5mm}
  \centering
    \includegraphics[width=0.57\textwidth]{images/trust_game_plot.pdf}
    \vspace{-7.5mm}
    \caption{\textbf{Amount Sent Distribution of LLM Agents and Humans as the Trustor in the Trust Game.} The size of  circles represents the number of personas for each amount sent.
The bold lines show the medians. 
The \textbf{crosses} indicate the \textbf{VRR} (\%) for different LLMs.}
  \label{fig:trust}
  \vspace{-4mm}
\end{wrapfigure}




\vspace{-3.5mm}
\section{Do LLM Agents Manifest Trust Behavior?}
\label{Do LLM Agents Manifest Trust Behavior?}
\vspace{-1.5mm}
In this section, we investigate whether or not LLM agents manifest trust behavior by letting LLM agents play the Trust Game (Section~\ref{game:Dictator Game and trust game} Game 1). 
In Behavioral Economics,  trust is widely measured by the initial amount sent from the trustor to the trustee in the Trust Game~\citep{Measuring_Trust,cesariniHeritabilityCooperativeBehavior2008}. Following the measurement of trust in human studies and the assumption humans own reasoning processes that underlie their decisions, we can define the conditions that LLM agents manifest trust behavior in the Trust Game as follows. 
\textit{First}, \textbf{the amount sent is positive and does not exceed the amount of money the trustor initially possesses}, which implies that the trustor places self-interest at risk with the expectation the trustee will reciprocate and that the trustor understands the money limit that can be given.  
\textit{Second}, \textbf{the decision (\ie, amounts sent) can be interpreted as the reasoning process (\ie, the BDI) of the trustor}. We explored utilizing BDI to model the reasoning process of LLM agents. If we can interpret the decision as the articulated reasoning process, we have evidence that LLM agents do not send a random amount of money and manifest some degree of rationality in the decision-making process.
Then, we assess whether LLM agents exhibit trust behavior based on two aspects:  the amount sent and the BDI.














\vspace{-2mm}
\subsection{Amount Sent}
\label{Amount Sent}
\vspace{-1mm}
To evaluate LLMs' capacity to understand the basic experimental setting regarding money limits, we propose a new evaluation metric, Valid Response Rate (VRR) (\%), defined as the percentage of personas with the amount sent falling within the initial money ($\$10$). Results are shown in Figure~\ref{fig:trust}. We can observe that \textbf{most LLMs have a high VRR except Llama-7b}, which implies that most LLMs manifest a full understanding regarding limits on the amount they can send in the Trust Game. Then, we observe the distribution of amounts sent for different LLMs as the trustor agent and discover that \textbf{the amounts sent are predominantly positive, indicating a level of trust}. 





\vspace{-2mm}
\subsection{Belief-Desire-Intention (BDI)}
\vspace{-1mm}

The sole evidence of the amount sent cannot sufficiently support the existence of trust behavior, because agents could send positive but random amounts of money. Thus, we leveraged the Belief-Desire-Intention framework~\citep{raoBDIAgentsTheory1995a,andreas-2022-language-BDI} to model the reasoning process of LLM agents. If we can interpret the amounts sent from BDI outputs, we have evidence to refute the hypothesis that the amounts sent are positive but random and demonstrate that LLM agents manifest some degree of  rationality.  We take GPT-4 as an example to analyze its BDI outputs. More examples from the other nine LLMs such as Vicuna-v1.3-7b are shown in the Appendix~\ref{BDI Analysis}.
Considering that the amounts sent typically vary across distinct personas, we select one BDI from the personas that give a high amount of money and another BDI from those  that give a low amount. \colorlightblue{Positive} and \colorlightred{negative} factors for trust behavior in the reasoning process are marked in blue and red, respectively.

\vspace{-2mm}
\begin{quoting}[leftmargin=10pt, rightmargin=10pt]
\textit{As a person with a  \colorlightblue{strong belief in the goodness of humanity}, \colorlightblue{I trust} that the other player ...Therefore, my desire is to maximize the outcome for both of us and \colorlightblue{cement a sense of comradery and trust}... I intend to use this as an opportunity to \colorlightblue{add what I can to someone else's life}...Finally, I will give \textbf{\underline{10 dollars}}.}
\end{quoting}
\vspace{-2mm}
We can observe that this persona shows a high-level of ``\colorlightblue{comradery and trust}'' towards the other player, which justifies the high amount sent from this persona (\ie, \textbf{\textit{\underline{10 dollars}}}).

\vspace{-2mm}
\begin{quoting}[leftmargin=10pt, rightmargin=10pt]
\textit{As an Analyst,.... My desire is that the other player will also see the benefits of \colorlightblue{reciprocity and goodwill} ... my \colorlightblue{intention is to give away a significant portion} of my initial 10 ... \colorlightred{However, since I have no knowledge of the other player}, ... Therefore, I \colorlightred{aim to give an amount that is not too high}, ...Finally, I will give \textbf{\underline{5 dollars}} to the other player...}
\end{quoting}
\vspace{-2mm}
Compared to the first persona, we see that the second one has a more cautious attitude. For example, ``\colorlightred{since I have no knowledge of the other player}'' shows skepticism regarding the other player's motives. Thus, this persona, though still optimistic about the other player (``\colorlightblue{intention ... give away a significant portion}''), strategically balances risk and  reciprocity, and then decides to send only a modest amount. 

Based on  GPT-4's BDI examples and examples from other LLMs in Appendix~\ref{BDI Analysis}, we find \textbf{decisions (\ie, amounts sent) from LLM agents in the Trust Game can be interpreted from their articulated reasoning process (\ie,  BDI)}.
Because most LLM agents have a high VRR--send a positive amount of money--and show some degree of rationality in giving money,
our first core finding is:








\vspace{-2mm}
\begin{center}
\begin{tcolorbox}[width=0.99\linewidth, boxrule=3pt, colback=gray!20, colframe=gray!20]
\textbf{Finding 1:} LLM agents generally exhibit trust behavior under the framework of the Trust Game.
\end{tcolorbox}
\end{center}
\vspace{-4mm}
\subsection{Basic Analysis of Agent Trust}
\vspace{-1.5mm}
We also conduct a basic analysis of LLM agents' trust behavior, namely agent trust, based on the results in Figure~\ref{fig:trust}. \textit{First}, we observe that Vicuna-7b has the highest level of trust towards the other player and GPT-3.5-turbo-0613 has the lowest level of trust as trust can be measured by the amount sent in human studies~\citep{Measuring_Trust,cesariniHeritabilityCooperativeBehavior2008}. \textit{Second}, compared with humans' average amount sent ($\$5.97$), most personas for GPT-4 and Vicuna-7b send a higher amount of money to the other player, and most personas for LLMs such as GPT-3.5-turb-0613 send a lower amount. \textit{Third}, we see that amounts sent for Llama2-70b and Llama2-13b have a convergent distribution while amounts sent for humans and Vicuna-7b are more divergent.


















